%=====================================================================
%  SECCION 3.4 -- ANALISIS TECNICO, ETICO, SOCIAL, AMBIENTAL
%  Responsable: Fabian Moreno Ugarte
%
%  GENERADO desde secciones/03_analisis_impacto.md, que es su FUENTE.
%  No editar aqui: editar el .md y reconvertir.
%
%  Resumen: el desarrollo completo esta en el Anexo C.
%=====================================================================


% ---- 3.4 ----
\subsection{Análisis Técnico, Ético, Social, Ambiental}
\label{sec:analisis}
\label{sec:eipd}
\label{sec:sintesis-impacto}

Este análisis sigue la estructura de una Evaluación de Impacto en
Protección de Datos: identificar el tratamiento, valorar su necesidad y
proporcionalidad, identificar los riesgos para los titulares y proponer
medidas de mitigación. Esa evaluación es exigible antes de iniciar un
tratamiento basado en la captación masiva y continua de imágenes, con
independencia de que el modelo produzca conteos agregados, de modo que
este apartado es su contenido sustantivo. Se prefiere declarar un
impacto como no cuantificado antes que estimarlo sin base. El desarrollo
completo figura en el Anexo C.

% ---- 3.4.1 ----
\subsubsection{Análisis Técnico}

El riesgo técnico dominante es la brecha de dominio. Los autores de P2R
\cite{Lin_2025_CVPR} reconocen que el desempeño de los métodos de conteo
se degrada de forma marcada cuando el entorno de despliegue difiere del
de entrenamiento, y documentan un modo de fallo que produce
sobreestimación severa del conteo. Ninguna cifra obtenida sobre datos
públicos predice, por tanto, la exactitud en el terminal. A ello se
suman el sesgo de anotación, pues la métrica de localización más
exigente de P2PNet \cite{Song_2021_ICCV} cae por efecto de las
desviaciones de etiquetado, y el sesgo demográfico, que ninguno de los
trabajos revisados evalúa de forma desagregada: la estatura, los
cubrimientos de cabeza, el equipaje voluminoso o las sillas de ruedas
pueden inducir desempeño desigual en una población tan heterogénea como
la de un aeropuerto internacional. Se recomienda que la validación en
sitio incluya una evaluación desagregada, al menos por franja de
estatura y por presencia de equipaje voluminoso.

% ---- 3.4.2 ----
\subsubsection{Análisis Ético}
\label{sec:impacto-etico}

El sistema ofrece una alternativa menos invasiva que la identificación
biométrica para la misma finalidad, pero su riesgo principal no es
técnico sino de gobernanza: el deslizamiento de finalidad. Una vez
desplegada la infraestructura, incorporar reidentificación para medir
tiempos de permanencia, o reconocimiento facial para detectar personas
de interés, es una decisión de configuración y de presupuesto, no un
rediseño. La mitigación que se propone es documental y contractual: una
declaración de finalidades excluidas conforme a la cual el sistema no
identifica personas, no realiza seguimiento individual de trayectorias,
no reconoce emociones y no alimenta decisiones automatizadas sobre
pasajeros individuales, reforzada por los requisitos de no persistencia
y salida agregada de la Sección 5. Una lista de lo que el sistema no
hace es más difícil de erosionar silenciosamente que una lista de lo que
sí hace.

% ---- 3.4.3 ----
\subsubsection{Análisis Social}
\label{sec:impacto-social}

El beneficio social es la prevención de incidentes por aglomeración y
una asignación de recursos basada en evidencia. Frente a él se
identifican tres riesgos. Si el subconteo se concentra en un grupo, las
zonas donde ese grupo se concentra reciben sistemáticamente menos
personal y menos apertura de mostradores: el daño no es individual sino
distributivo. El segundo es el sesgo de automatización: si el personal
llega a confiar en el tablero más que en su propia observación, un fallo
en un momento crítico produce un efecto peor que la ausencia del
sistema, y un operador que aprueba mecánicamente sus alertas reintroduce
de hecho la decisión automatizada que la supervisión humana pretende
evitar. Por eso se recomienda tratar la capacitación de los operadores
en los modos de fallo del modelo como medida de mitigación y no como
accesorio de implantación. El tercero es el uso laboral no declarado: si
los conteos por zona se cruzan con los turnos del personal, el sistema
se convierte en una herramienta de evaluación de desempeño, y esa
exclusión debe figurar expresamente en la declaración de finalidades.

% ---- 3.4.4 ----
\subsubsection{Análisis Ambiental}
\label{sec:impacto-ambiental}

El equipo no ha medido el consumo energético del entrenamiento ni de la
inferencia, de modo que este análisis es cualitativo. La reutilización
de modelos preentrenados evita entrenar desde cero, y la pérdida de P2R
prescinde del algoritmo húngaro: sobre una imagen de 576 × 960 píxeles
con 775 puntos anotados, su cálculo es cerca de sesenta y ocho veces más
rápido que el de la pérdida punto a punto \cite{Lin_2025_CVPR}, aunque
esa cifra mide una operación por iteración y no el consumo total. A lo
largo de la vida útil, la inferencia continua pesará previsiblemente más
que el entrenamiento, y las decisiones que la reducen (muestrear un
fotograma cada varios segundos, procesar en el borde, no persistir
imágenes y agregar la salida) son las mismas que reducen la exposición
de datos personales. Se recomienda registrar las horas de GPU desde el
primer entrenamiento, porque ese dato no es reconstruible a posteriori,
y reutilizar la infraestructura de CCTV existente, que es además la
opción económicamente decisiva.
